% ================= BÌA =================
\begin{titlepage}
\begin{tikzpicture}[remember picture,overlay]
  \node[inner sep=0] at (current page.center)
    {\includegraphics[width=\paperwidth,height=\paperheight]{cover-bg.jpg}};
  \node[font=\bfseries\fontsize{16}{20}\selectfont] at ([yshift=-45mm]current page.north) {\tenkhoa};
  \node[font=\bfseries,align=center] at ([yshift=-123mm]current page.north)
    {BÁO CÁO KIẾN TẬP\\[4pt] NGÀNH \tennganh\\[4pt] (CHUYÊN NGÀNH \tenchuyennganh)};
  \node[font=\bfseries,align=center,text width=16.5cm] at ([yshift=-164mm]current page.north) {\tendetai};
  \node[anchor=north west,align=left] at ([xshift=90mm,yshift=-194mm]current page.north west)
    {\textbf{GVHD:} \gvhd\\[4pt] \textbf{SVTH:} \svth\\[4pt] \textbf{MSSV:} \mssv};
  \node[font=\bfseries,text=white] at ([yshift=-271mm]current page.north) {\thoigian};
\end{tikzpicture}
\end{titlepage}
\setcounter{page}{1}

% ================= PHỤ BÌA =================
\clearpage
\todo{Trang phụ bìa theo mẫu UEL (giống bìa, không ảnh nền).}

% ================= LỜI CẢM ƠN =================
\heading{LỜI CẢM ƠN}
\todo{Cảm ơn giảng viên hướng dẫn, Mentor tại cơ quan kiến tập; mong nhận góp ý.}

% ================= LỜI CAM KẾT =================
\heading{LỜI CAM KẾT}
\todo{Cam kết nội dung trung thực, do bản thân thực hiện, trích dẫn đầy đủ.}

% ================= PHIẾU ĐÁNH GIÁ KẾT QUẢ KIẾN TẬP =================
\clearpage
\section*{PHIẾU ĐÁNH GIÁ KẾT QUẢ KIẾN TẬP}
{\setlength{\parskip}{0pt}\singlespacing
MSSV: \mssv

Họ tên: \svth

Tên Cơ quan Kiến tập: \coquan

Địa chỉ Cơ quan: \dotfill

Điện thoại Cơ quan: \dotfill

Họ tên người đại diện cơ quan: \mentor

Chức vụ: \makebox[5cm]{\dotfill} Điện thoại: \dotfill

Thời gian kiến tập tại Cơ quan: Từ ngày ../../.... Đến ngày: ../../....

\textit{* Đề nghị \textbf{Quý Cơ quan} đánh giá \textbf{bằng cách đánh dấu X} vào cột xếp loại các nội dung đánh giá trong bảng sau.}

Ghi chú:

\textbf{\textit{Loại A: 2đ; Loại B: 1,5đ}}

\textbf{\textit{Loại C: 1đ; Loại D: 0,5đ}}
\par}
\vspace{6pt}
{\renewcommand{\arraystretch}{1.4}\noindent
\begin{tabularx}{\textwidth}{|X|c|c|c|c|}\hline
\multirow{2}{*}{\centering\textbf{Nội dung đánh giá}} & \multicolumn{4}{c|}{\textbf{Xếp loại}} \\ \cline{2-5}
 & \makebox[0.8cm]{A} & \makebox[0.8cm]{B} & \makebox[0.8cm]{C} & \makebox[0.8cm]{D} \\ \hline
\multicolumn{5}{|l|}{\textbf{1. Tinh thần kỷ luật, thái độ}} \\ \hline
\quad 1.1 Thực hiện nội quy của cơ quan & & & & \\ \hline
\quad 1.2 Chấp hành giờ giấc làm việc & & & & \\ \hline
\quad 1.3 Thái độ giao tiếp với đồng nghiệp, cộng sự & & & & \\ \hline
\quad 1.4 Ý thức bảo vệ tài sản & & & & \\ \hline
\quad 1.5 Tích cực trong công việc & & & & \\ \hline
\multicolumn{5}{|l|}{\textbf{2. Kỹ năng chuyên môn, nghiệp vụ}} \\ \hline
\quad 2.1 Đáp ứng yêu cầu công việc & & & & \\ \hline
\quad 2.2 Tinh thần học hỏi, nâng cao trình độ chuyên môn & & & & \\ \hline
\quad 2.3 Có sáng kiến, năng động trong công việc & & & & \\ \hline
\multicolumn{5}{|l|}{\textbf{3. Kết quả thực tập}} \\ \hline
\quad 3.1 Có sản phẩm ứng dụng thực tế đem lại lợi ích cho Cơ quan & & & & \\ \hline
\quad 3.2 Mức độ hoàn thành nhiệm vụ thực tập & & & & \\ \hline
\end{tabularx}}

% ================= PHIẾU NHẬN XÉT CƠ QUAN KIẾN TẬP =================
\clearpage
\section*{PHIẾU NHẬN XÉT CƠ QUAN KIẾN TẬP}
\dongcham{12}
\vspace{12pt}
\hfill\begin{minipage}{8cm}\centering
............, ngày ..... tháng ..... năm .......\\
\textbf{XÁC NHẬN CỦA CƠ QUAN}\\
\textit{Ký, ghi rõ họ tên, đóng dấu}
\end{minipage}

% ================= PHIẾU ĐÁNH GIÁ CỦA GVHD =================
\clearpage
\section*{PHIẾU ĐÁNH GIÁ BÁO CÁO KIẾN TẬP\\ CỦA GIẢNG VIÊN HƯỚNG DẪN}
{\setlength{\parskip}{0pt}\singlespacing
MSSV: \mssv

Họ tên: \svth

Lớp: \lop

Ngành/Chuyên ngành: \tennganh{} (Chuyên ngành \tenchuyennganh)

Tên đề tài: \tendetai

GVHD: \gvhd
\par}
\vspace{6pt}
{\singlespacing\renewcommand{\arraystretch}{1.3}\noindent
\begin{tabularx}{\textwidth}{|c|>{\centering\arraybackslash}X|c|>{\centering\arraybackslash}p{2.2cm}|}\hline
\textbf{STT} & \textbf{Tiêu chí cụ thể} & \textbf{Tỉ lệ} & \textbf{Điểm đã nhân tỉ lệ} \\ \hline
1 & Tuân thủ quy định học phần, đúng tiến độ, đầy đủ xác nhận của đơn vị/leader, tinh thần học tập nghiêm túc & 10\% & \\ \hline
2 & Tìm hiểu doanh nghiệp, môi trường làm việc, vị trí nghề nghiệp, xác định định hướng chuyên sâu rõ ràng & 10\% & \\ \hline
3 & Mô tả bối cảnh thực tế, xác định vấn đề/tình huống rõ ràng, có logic và phù hợp với định hướng nghề nghiệp & 15\% & \\ \hline
4 & Vận dụng kiến thức chuyên ngành, kiến thức kinh doanh quản lý, công nghệ để phân tích vấn đề & 20\% & \\ \hline
5 & Thu thập dữ liệu, xử lý dữ liệu, phân tích, diễn giải kết quả, sử dụng bảng biểu/hình minh họa phù hợp & 15\% & \\ \hline
6 & Đề xuất giải pháp phù hợp, có tính ứng dụng, khả thi trong bối cảnh thực tế & 15\% & \\ \hline
7 & Bố cục logic, trình bày khoa học, ngôn ngữ học thuật, trích dẫn tài liệu, hình thức đúng quy định & 15\% & \\ \hline
\multicolumn{3}{|c|}{Tổng điểm:} & \\[12pt] \hline
\end{tabularx}}

\vspace{6pt}
\hfill\begin{minipage}{8cm}\centering
............, ngày ..... tháng ..... năm .......\\
\textbf{Giảng viên hướng dẫn}\\
\textit{Ký, ghi rõ họ tên}
\end{minipage}

% ================= PHIẾU NHẬN XÉT CỦA GVHD =================
\clearpage
\section*{PHIẾU NHẬN XÉT CỦA GIẢNG VIÊN HƯỚNG DẪN}
\dongcham{19}
\vspace{6pt}
\hfill\begin{minipage}{8cm}\centering
............, ngày ..... tháng ..... năm .......\\
\textbf{Giảng viên hướng dẫn}\\
\textit{Ký, ghi rõ họ tên}
\end{minipage}

% ================= LỊCH LÀM VIỆC =================
\heading{LỊCH LÀM VIỆC VÀ NỘI DUNG CÔNG VIỆC}
{\singlespacing\renewcommand{\arraystretch}{1.4}\noindent
\begin{tabularx}{\textwidth}{|c|p{2.6cm}|X|p{3.6cm}|}\hline
\textbf{Tuần} & \textbf{Thời gian} & \textbf{Nội dung công việc} & \textbf{Sản phẩm / minh chứng} \\ \hline
1 & ../../ -- ../../ & Tìm hiểu MCP, agentgateway, Kubernetes/kind & Môi trường kind \\ \hline
2 & ../../ -- ../../ & Xây dựng mock MCP server (FastMCP, mockkit) & \texttt{servers/} \\ \hline
3 & ../../ -- ../../ & Chuẩn hóa danh mục 21 tool và ma trận quyền & \texttt{agent-tool-access.json} \\ \hline
4 & ../../ -- ../../ & Thiết lập AI Agent trên Keycloak & \texttt{setup-agents.sh} \\ \hline
5 & ../../ -- ../../ & Observability: access log, Loki, Grafana & Dashboard allow/deny \\ \hline
6 & ../../ -- ../../ & Tự động hóa \texttt{ops/up.ps1}, \texttt{ops/down.ps1} & Log chạy script \\ \hline
7 & ../../ -- ../../ & Kiểm chứng, thu thập bằng chứng, viết báo cáo & \texttt{evidence/}, báo cáo \\ \hline
\end{tabularx}}
\todo{Điều chỉnh theo thời gian và công việc thực tế (đối chiếu git log).}

% ================= MỤC LỤC & CÁC DANH MỤC =================
\clearpage
\addcontentsline{toc}{section}{MỤC LỤC}
{\setlength{\parskip}{0pt}\tableofcontents}

\clearpage
\addcontentsline{toc}{section}{DANH MỤC HÌNH ẢNH}
{\setlength{\parskip}{0pt}\renewcommand{\numberline}[1]{Hình #1: }\listoffigures}

\clearpage
\addcontentsline{toc}{section}{DANH MỤC BẢNG}
{\setlength{\parskip}{0pt}\renewcommand{\numberline}[1]{Bảng #1: }\listoftables}

\heading{DANH MỤC TỪ VIẾT TẮT}
{\singlespacing\renewcommand{\arraystretch}{1.5}
\begin{longtable}{|>{\centering\arraybackslash}p{2.2cm}|>{\centering\arraybackslash}p{6.1cm}|>{\centering\arraybackslash}p{6.1cm}|}\hline
\textbf{Từ viết tắt} & \textbf{Cụm từ đầy đủ} & \textbf{Diễn giải tiếng Việt} \\ \hline
\endhead
AI & Artificial Intelligence & Trí tuệ nhân tạo \\ \hline
API & Application Programming Interface & Giao diện lập trình ứng dụng \\ \hline
CEL & Common Expression Language & Ngôn ngữ biểu thức của chính sách \\ \hline
CRD & Custom Resource Definition & Tài nguyên tùy biến của Kubernetes \\ \hline
JWT & JSON Web Token & Mã thông báo dạng JSON \\ \hline
JWKS & JSON Web Key Set & Tập khóa công khai xác minh JWT \\ \hline
K8s & Kubernetes & Nền tảng điều phối container \\ \hline
MCP & Model Context Protocol & Giao thức ngữ cảnh mô hình \\ \hline
OIDC & OpenID Connect & Giao thức xác thực trên OAuth 2.0 \\ \hline
PoC & Proof of Concept & Bản chứng minh khái niệm \\ \hline
RBAC & Role-Based Access Control & Kiểm soát truy cập theo vai trò \\ \hline
% ... bổ sung các từ viết tắt khác
\end{longtable}}

\heading{DANH MỤC THUẬT NGỮ}
{\singlespacing\renewcommand{\arraystretch}{1.5}
\begin{longtable}{|>{\centering\arraybackslash}p{1.5cm}|>{\centering\arraybackslash}p{6.2cm}|>{\centering\arraybackslash}p{6.9cm}|}\hline
\textbf{STT} & \textbf{Thuật ngữ gốc} & \textbf{Ý nghĩa} \\ \hline
\endhead
1 & AI Agent & Tác tử AI tự gọi công cụ để hoàn thành tác vụ \\ \hline
2 & MCP Gateway & Cổng tập trung xác thực, phân quyền, định tuyến MCP \\ \hline
3 & Tool & Hàm mà MCP server cung cấp cho tác tử \\ \hline
4 & Mock server & Máy chủ giả lập trả dữ liệu mẫu \\ \hline
5 & Deny-by-default & Mặc định từ chối, chỉ cho phép khi có quy tắc \\ \hline
6 & Observability & Khả năng quan sát qua log, metric, trace \\ \hline
7 & Access log & Nhật ký truy cập \\ \hline
% ... bổ sung các thuật ngữ khác
\end{longtable}}
